% GENERATED FILE. Edit canonical lessons, then run npm run generate:books.
% canonical-source-hash: fnv1a64:7da11674ceb7013e
% canonical-lessons: LA-C309-accuso, LA-C309-accusator, LA-C309-advocatus, LA-C309-patronus, LA-C309-testimonium

\chapter{Accūsātor et patrōnus}
\label{ch:la-accusator-et-patronus}
\hlchaptermodality{309}

\begin{chapteropening}
\textbf{By the end of this chapter:} \emph{I can describe the sides in a trial: to accuse, the prosecutor, the supporter in court, the counsel and the evidence.}

Complete the last lesson of chapter 309: I can describe the sides in a trial: to accuse, the prosecutor, the supporter in court, the counsel and the evidence.
\end{chapteropening}

\section[accūsō, accūsāre]{accūsō, accūsāre — to accuse, to charge}

\label{lesson:LA-C309-accuso}



\begin{quote}
Before the new one: say the Latin for a law, then the Latin for law.
\end{quote}

\subsection*{You'll want to know: accūsō, accūsāre}
\textbf{accūsō, accūsāre} — \textquotedblleft{}to accuse, to charge\textquotedblright{}.

\begin{quote}
Mercātōrem fūrtī accūsāvērunt.
\end{quote}

\begin{quote}
\emph{They accused the merchant of theft.}
\end{quote}

\emph{accūsō, accūsāre, accūsāvī, accūsātum}, with the genitive of the charge: \emph{fūrtī accūsāre}.

\begin{cousinweb}
\textbf{accūsō} is built on \emph{causa}, \textquotedblleft{}a cause, a lawsuit\textquotedblright{}, with \emph{ad-} (EtymDB): to call someone to account. English \emph{accuse} comes from it.
\end{cousinweb}

\begin{grammarlens}[title={What you've built}]
One more word for telling, explaining and describing at B1.
\end{grammarlens}

\subsection*{Your turn}
\begin{itemize}
\raggedright
  \item \emph{Say it:} \emph{accūsō, accūsāre}
  \item \emph{Say it:} \emph{accūsō, accūsāre}, once more
  \item \emph{Say it:} a sentence of your own with \emph{accūsō, accūsāre}
  \item \emph{Recall:} say the Latin for a law, then the Latin for law, then say \emph{accūsō, accūsāre} again
\end{itemize}

\begin{tcolorbox}[breakable,colback=teal!4,colframe=teal!35!black,title={Before you move on}]
What does accūsō, accūsāre mean? (\textquotedblleft{}To accuse, to charge\textquotedblright{}.) Say it once more.
\end{tcolorbox}
\section[accūsātor, accūsātōris]{accūsātor, accūsātōris — a prosecutor, an accuser}

\label{lesson:LA-C309-accusator}



\begin{quote}
Before the new one: say the Latin for law, then the Latin for to accuse.
\end{quote}

\subsection*{You'll want to know: accūsātor, accūsātōris}
\textbf{accūsātor, accūsātōris} — \textquotedblleft{}a prosecutor, an accuser\textquotedblright{}.

\begin{quote}
Accūsātor longam ōrātiōnem habuit.
\end{quote}

\begin{quote}
\emph{The prosecutor made a long speech.}
\end{quote}

Masculine, 3rd declension. Rome had no public prosecutor: any citizen could act as \emph{accūsātor}, and young men won fame that way.

\begin{cousinweb}
\textbf{accūsātor} is built on \emph{accūsō}, \textquotedblleft{}to accuse\textquotedblright{} (EtymDB), with the ending \emph{-tor} for the one who does it.
\end{cousinweb}

\begin{grammarlens}[title={What you've built}]
One more word for telling, explaining and describing at B1.
\end{grammarlens}

\subsection*{Your turn}
\begin{itemize}
\raggedright
  \item \emph{Say it:} \emph{accūsātor, accūsātōris}
  \item \emph{Say it:} \emph{accūsātor, accūsātōris}, once more
  \item \emph{Say it:} a sentence of your own with \emph{accūsātor, accūsātōris}
  \item \emph{Recall:} say the Latin for law, then the Latin for to accuse, then say \emph{accūsātor, accūsātōris} again
\end{itemize}

\begin{tcolorbox}[breakable,colback=teal!4,colframe=teal!35!black,title={Before you move on}]
What does accūsātor, accūsātōris mean? (\textquotedblleft{}A prosecutor, an accuser\textquotedblright{}.) Say it once more.
\end{tcolorbox}
\section[advocātus, advocātī]{advocātus, advocātī — a supporter in court, a legal adviser}

\label{lesson:LA-C309-advocatus}



\begin{quote}
Before the new one: say the Latin for to accuse, then the Latin for a prosecutor.
\end{quote}

\subsection*{You'll want to know: advocātus, advocātī}
\textbf{advocātus, advocātī} — \textquotedblleft{}a supporter in court, a legal adviser\textquotedblright{}.

\begin{quote}
Reus multōs amīcōs advocātōs habēbat.
\end{quote}

\begin{quote}
\emph{The defendant had many friends supporting him in court.}
\end{quote}

Masculine. In Cicero an \emph{advocātus} stood by a friend in court; later the word meant a pleader. English \emph{advocate}.

\begin{cousinweb}
\textbf{advocātus} is the participle of \emph{advocō}, \textquotedblleft{}to call to oneself\textquotedblright{} (EtymDB), from \emph{ad} and \emph{vocō}, \textquotedblleft{}to call\textquotedblright{}: one who has been called in. English \emph{advocate} comes from it.
\end{cousinweb}

\begin{grammarlens}[title={What you've built}]
One more word for telling, explaining and describing at B1.
\end{grammarlens}

\subsection*{Your turn}
\begin{itemize}
\raggedright
  \item \emph{Say it:} \emph{advocātus, advocātī}
  \item \emph{Say it:} \emph{advocātus, advocātī}, once more
  \item \emph{Say it:} a sentence of your own with \emph{advocātus, advocātī}
  \item \emph{Recall:} say the Latin for to accuse, then the Latin for a prosecutor, then say \emph{advocātus, advocātī} again
\end{itemize}

\begin{tcolorbox}[breakable,colback=teal!4,colframe=teal!35!black,title={Before you move on}]
What does advocātus, advocātī mean? (\textquotedblleft{}A supporter in court, a legal adviser\textquotedblright{}.) Say it once more.
\end{tcolorbox}
\section[patrōnus, patrōnī]{patrōnus, patrōnī — a patron; a defending counsel}

\label{lesson:LA-C309-patronus}



\begin{quote}
Before the new one: say the Latin for a prosecutor, then the Latin for a supporter in court.
\end{quote}

\subsection*{You'll want to know: patrōnus, patrōnī}
\textbf{patrōnus, patrōnī} — \textquotedblleft{}a patron; a defending counsel\textquotedblright{}.

\begin{quote}
Patrōnus prō reō dīxit et iūdicēs mōvit.
\end{quote}

\begin{quote}
\emph{The counsel spoke for the defendant and moved the judges.}
\end{quote}

Masculine. A patron protected his dependants and spoke for them in court. English \emph{patron}.

\begin{cousinweb}
\textbf{patrōnus} is built on \emph{pater}, \textquotedblleft{}father\textquotedblright{} (EtymDB): one who stands to others as a father. English \emph{patron} comes from it.
\end{cousinweb}

\begin{grammarlens}[title={What you've built}]
One more word for telling, explaining and describing at B1.
\end{grammarlens}

\subsection*{Your turn}
\begin{itemize}
\raggedright
  \item \emph{Say it:} \emph{patrōnus, patrōnī}
  \item \emph{Say it:} \emph{patrōnus, patrōnī}, once more
  \item \emph{Say it:} a sentence of your own with \emph{patrōnus, patrōnī}
  \item \emph{Recall:} say the Latin for a prosecutor, then the Latin for a supporter in court, then say \emph{patrōnus, patrōnī} again
\end{itemize}

\begin{tcolorbox}[breakable,colback=teal!4,colframe=teal!35!black,title={Before you move on}]
What does patrōnus, patrōnī mean? (\textquotedblleft{}A patron; a defending counsel\textquotedblright{}.) Say it once more.
\end{tcolorbox}
\section[testimōnium, testimōniī]{testimōnium, testimōniī — evidence, testimony}

\label{lesson:LA-C309-testimonium}



\begin{quote}
Before the new one: say the Latin for a supporter in court, then the Latin for a patron.
\end{quote}

\subsection*{You'll want to know: testimōnium, testimōniī}
\textbf{testimōnium, testimōniī} — \textquotedblleft{}evidence, testimony\textquotedblright{}.

\begin{quote}
Testimōnium ancillae ūtile erat.
\end{quote}

\begin{quote}
\emph{The maid's evidence was useful.}
\end{quote}

Neuter, from \emph{testis}, 'witness' (taught). \emph{Testimōnium dīcere}: to give evidence.

\begin{cousinweb}
\textbf{testimōnium} is built on \emph{testis}, \textquotedblleft{}witness\textquotedblright{} (EtymDB), with the ending \emph{-mōnium}. English \emph{testimony} comes from it.
\end{cousinweb}

\begin{grammarlens}[title={What you've built}]
One more word for telling, explaining and describing at B1.
\end{grammarlens}

\subsection*{Your turn}
\begin{itemize}
\raggedright
  \item \emph{Say it:} \emph{testimōnium, testimōniī}
  \item \emph{Say it:} \emph{testimōnium, testimōniī}, once more
  \item \emph{Say it:} a sentence of your own with \emph{testimōnium, testimōniī}
  \item \emph{Recall:} say the Latin for a supporter in court, then the Latin for a patron, then say \emph{testimōnium, testimōniī} again
\end{itemize}

\begin{tcolorbox}[breakable,colback=teal!4,colframe=teal!35!black,title={Before you move on}]
What does testimōnium, testimōniī mean? (\textquotedblleft{}Evidence, testimony\textquotedblright{}.) Say it once more.
\end{tcolorbox}
